\clearpage
\section{Phân phối mẫu: từ một quan sát đến thống kê của cả mẫu}
\label{sec:sampling}

\subsection{Trung bình mẫu và hiệp phương sai mẫu}
Cho $\mathbf X_1,\ldots,\mathbf X_n$ là mẫu iid, mỗi quan sát có $p$ biến.
Định nghĩa
\begin{align}
 \Xbar&=\frac1n\sum_{j=1}^n\mathbf X_j,
 \label{eq:samplemean}\\
 S&=\frac1{n-1}\sum_{j=1}^n
 (\mathbf X_j-\Xbar)(\mathbf X_j-\Xbar)'.
 \label{eq:samplecov}
\end{align}
Mỗi phần tử $s_{ik}$ là hiệp phương sai mẫu của hai cột $i,k$.
Ký hiệu $S$ dùng cho cả thống kê ngẫu nhiên và ma trận giá trị đã tính;
ngữ cảnh sẽ phân biệt. $\Xbar$ ngẫu nhiên trước khi lấy mẫu, còn $\xbar$
là giá trị cụ thể sau khi lấy mẫu.

\begin{example}[Tính từ bảng bốn quan sát]
Với $\mathbf D=\begin{pmatrix}2&12\\8&9\\6&9\\8&10\end{pmatrix}$,
\[
 \xbar=\begin{pmatrix}6\\10\end{pmatrix},\qquad
 \mathbf D_{\mathrm{centered}}=
 \begin{pmatrix}-4&2\\2&-1\\0&-1\\2&0\end{pmatrix}.
\]
Đường chéo và hiệp phương sai lần lượt là
\[
 s_{11}=\frac{16+4+0+4}{3}=8,\quad
 s_{22}=\frac{4+1+1+0}{3}=2,
\]
\[
 s_{12}=\frac{(-4)(2)+(2)(-1)+(0)(-1)+(2)(0)}3=-\frac{10}3.
\]
Do đó
\begin{equation}
 S=\begin{pmatrix}8&-10/3\\-10/3&2\end{pmatrix},\qquad
 S^{-1}=\frac1{22}\begin{pmatrix}9&15\\15&36\end{pmatrix}.
 \label{eq:sample4}
\end{equation}
Ma trận tích $\mathbf D_{\mathrm{centered}}'
\mathbf D_{\mathrm{centered}}$ cho đúng $(n-1)S$.
\end{example}

\subsection{Phân phối của trung bình mẫu chuẩn}
\begin{result}
Nếu $\mathbf X_j\overset{\mathrm{iid}}\sim\N_p(\vmu,\Sigma)$, thì
\begin{equation}
 \Xbar\sim\N_p\left(\vmu,\frac{\Sigma}{n}\right).
 \label{eq:samplemeandist}
\end{equation}
\end{result}
\nguon{Chương 4, §4.4, công thức (4-23), tr.~174}

Kỳ vọng vẫn bằng $\vmu$, còn hiệp phương sai giảm $n$ lần.
Điều này do các quan sát độc lập: các hiệp phương sai giữa đối tượng khác
nhau bằng 0. Với một tổ hợp $\mathbf a'\Xbar$, phương sai là
$\mathbf a'\Sigma\mathbf a/n$.

\hieu{Nếu lấy lại nhiều mẫu, trung bình của các mẫu sẽ tụ quanh cùng một
vectơ trung bình tổng thể. Mẫu càng lớn, đám mây các trung bình mẫu càng nhỏ.
Các quan sát riêng lẻ vẫn có độ phân tán $\Sigma$; chúng không tự trở nên
ít biến thiên chỉ vì ta lấy thêm dữ liệu.}

Với $p=2$, $n=25$ và $\Sigma=\diag(4,9)$, độ lệch chuẩn của hai thành
phần trung bình mẫu là $2/5$ và $3/5$, còn độ lệch chuẩn của quan sát là 2
và 3. Cần giữ rõ ta đang xét loại đại lượng nào.

\subsection{Dạng toàn phương khi biết hiệp phương sai tổng thể}
Từ \eqref{eq:samplemeandist}, $\sqrt n(\Xbar-\vmu)\sim\N_p(0,\Sigma)$.
Áp dụng kết quả khoảng cách Mahalanobis cho vectơ này:
\begin{equation}
 n(\Xbar-\vmu)'\Sigma^{-1}(\Xbar-\vmu)\sim\chi^2_p.
 \label{eq:knowncovmean}
\end{equation}
Đây là kết quả \emph{chính xác} trong mẫu chuẩn khi dùng $\Sigma$ thật,
ngay cả ở mẫu nhỏ. Trong Bài 2.9 với $n=60,p=4$, trung bình mẫu có
hiệp phương sai $\Sigma/60$; dạng toàn phương của một quan sát và dạng
\eqref{eq:knowncovmean} đều có phân phối $\chi^2_4$.
Không có mâu thuẫn: hệ số $n$ bù cho việc trung bình mẫu biến thiên ít hơn.
\nguon{Hệ quả của Chương 4, Result 4.7, tr.~163, và (4-23), tr.~174}

Nếu thay $\Sigma^{-1}$ bằng $S^{-1}$, ma trận chuẩn hóa cũng ngẫu nhiên.
Khi đó không thể giữ nguyên kết luận chi-bình phương chính xác; cần
phân phối Hotelling hữu hạn hoặc xấp xỉ mẫu lớn. Đây là điểm phân biệt
quan trọng giữa các câu của Bài 2.9.

\subsection{Wishart: mở rộng chi-bình phương từ số sang ma trận}
Trong một chiều, tổng bình phương các biến chuẩn độc lập cho phân phối
chi-bình phương. Trong nhiều chiều, thay bình phương $Z^2$ bằng tích ngoài
$\mathbf Z\mathbf Z'$, ta có phân phối Wishart.

\begin{definition}
Nếu $\mathbf Z_1,\ldots,\mathbf Z_\nu$ iid $\N_p(0,\Sigma)$, thì
\[
 W=\sum_{j=1}^\nu\mathbf Z_j\mathbf Z_j'
 \sim W_{p,\nu}(\Sigma).
\]
$p$ là kích thước ma trận, $\nu$ là số bậc tự do, $\Sigma$ là ma trận
tham số tỉ lệ. Một số tài liệu viết $W_p(\Sigma,\nu)$ cho cùng phân phối.
\end{definition}
\nguon{Chương 4, §4.4, định nghĩa (4-22), tr.~174}

\begin{result}
Với mẫu iid chuẩn nhiều chiều,
\begin{equation}
 (n-1)S\sim W_{p,n-1}(\Sigma),\qquad \Xbar\perp S.
 \label{eq:wishart}
\end{equation}
\end{result}
\nguon{Chương 4, §4.4, công thức (4-23), tr.~174}

\hieu{Trung bình mẫu cung cấp thông tin về vị trí; hiệp phương sai mẫu
cung cấp thông tin về độ phân tán. Trong mô hình chuẩn, hai phần thông tin
này ngẫu nhiên độc lập với nhau. Tính chất đặc biệt đó cho phép xây dựng
phân phối kiểm định chính xác ở phần tiếp theo.}

Vì sao bậc tự do là $n-1$? Các độ lệch cộng lại bằng 0:
$\sum_j(\mathbf X_j-\Xbar)=0$. Sau khi biết $n-1$ độ lệch, độ lệch còn
lại đã được xác định. Các độ lệch này không iid, nhưng có thể đổi hệ tọa độ
giữa các quan sát để viết tổng tích ngoài thành tổng của $n-1$ vectơ chuẩn
độc lập. Phụ lục~\ref{app:wishart} mô tả phép đổi đó.

Với $p=1$, Wishart thu về
\[
 \frac{(n-1)s^2}{\sigma^2}\sim\chi^2_{n-1},
\]
nên đây thực sự là mở rộng của kết quả một chiều.

\subsection{Khi nào \texorpdfstring{$S$}{S} có nghịch đảo?}
Hạng của ma trận dữ liệu đã trừ trung bình không vượt quá $n-1$. Do đó
\[
 \operatorname{rank}(S)\le\min(p,n-1).
\]
Nếu $n\le p$, $S$ chắc chắn suy biến. Nếu mẫu iid chuẩn, $\Sigma$ xác định
dương và $n>p$, thì $S$ khả nghịch với xác suất 1. Trong dữ liệu thực,
biến lặp lại hoặc quan hệ tuyến tính chính xác vẫn có thể gây suy biến.

\luuy{Điều kiện $n>p$ là cần thiết để dùng nghịch đảo hiệp phương sai
mẫu trong kiểm định cổ điển. Nó không tự bảo đảm dữ liệu bất kỳ không có
cột trùng hay cột hằng. Cũng không nên thay $S^{-1}$ bằng nghịch đảo giả rồi
giữ nguyên ngưỡng $F$; đó là một phương pháp khác.}

\subsection{Thiết lập MLE: đi từ mật độ đến ước lượng}
\label{sec:mle_derivation}
Ước lượng hợp lý cực đại (maximum likelihood estimation, MLE) chọn các
tham số làm dữ liệu \emph{đã quan sát} có mật độ chung lớn nhất. Khi tìm
MLE, các $\mathbf x_j$ là số đã biết; $\vmu,\Sigma$ là các biến cần tối ưu.
Ta không tìm tham số có xác suất lớn nhất, vì tham số ở đây không được
đặt một phân phối xác suất.

Giả sử mẫu iid $\N_p(\vmu,\Sigma)$ và $\Sigma\succ0$. Để có nghiệm MLE
không suy biến, giả sử ma trận tổng tích ngoài quanh $\xbar$ xác định dương.
Trong mẫu chuẩn với hiệp phương sai tổng thể xác định dương, $n>p$ bảo
đảm điều này với xác suất 1; đối với bảng số cụ thể vẫn phải kiểm tra hạng.

\buoc{1}{Viết hàm hợp lý và log-hợp lý}
Tính độc lập cho phép nhân $n$ mật độ:
\begin{equation}
 L(\vmu,\Sigma)=(2\pi)^{-np/2}|\Sigma|^{-n/2}
 \exp\!\left[-\frac12\sum_{j=1}^n
 (\mathbf x_j-\vmu)'\Sigma^{-1}(\mathbf x_j-\vmu)\right].
 \label{eq:mlelikelihood}
\end{equation}
Logarit là hàm tăng, nên cực đại $L$ cũng là cực đại $\ell=\log L$:
\begin{equation}
 \ell(\vmu,\Sigma)=-\frac{np}2\log(2\pi)-\frac n2\log|\Sigma|
 -\frac12\sum_{j=1}^n(\mathbf x_j-\vmu)'\Sigma^{-1}(\mathbf x_j-\vmu).
 \label{eq:mleloglikelihood}
\end{equation}
Logarit chuyển tích thành tổng, làm phép tối ưu dễ xử lý hơn.
\nguon{Chương 4, §4.3, mật độ mẫu (4-11), tr.~168; hàm hợp lý (4-16), tr.~170}

\buoc{2}{Tách phần phụ thuộc trung bình và tìm $\widehat{\vmu}$}
Đặt
\[
 A=\sum_j(\mathbf x_j-\xbar)(\mathbf x_j-\xbar)'=(n-1)S.
\]
Vì $\sum_j(\mathbf x_j-\xbar)=0$, hai hạng chéo triệt tiêu khi khai triển:
\begin{equation}
 \sum_j(\mathbf x_j-\vmu)(\mathbf x_j-\vmu)'
 =A+n(\xbar-\vmu)(\xbar-\vmu)'.
 \label{eq:mlescatterdecomposition}
\end{equation}
Trace, ký hiệu $\tr(B)$, là tổng các phần tử đường chéo của ma trận vuông.
Với vectơ $\mathbf u$, ta có
$\mathbf u'B\mathbf u=\tr(B\mathbf u\mathbf u')$; do đó
\begin{align}
 \ell(\vmu,\Sigma)
 &=K-\frac n2\log|\Sigma|-\frac12\tr(\Sigma^{-1}A)
 \notag\\[-2pt]
 &\qquad-\frac n2(\xbar-\vmu)'\Sigma^{-1}(\xbar-\vmu),
 \label{eq:mlemeansplit}
\end{align}
trong đó $K=-np\log(2\pi)/2$ không phụ thuộc tham số.

Với $\Sigma$ cố định, chỉ hạng cuối phụ thuộc $\vmu$. Vì
$\Sigma^{-1}\succ0$, hạng toàn phương không âm, bằng 0 duy nhất khi
$\vmu=\xbar$. Nó mang dấu trừ trong log-hợp lý, nên
\begin{equation}
 \widehat{\vmu}=\xbar.
 \label{eq:mlemeanhat}
\end{equation}
Kết quả đúng với \emph{mọi} $\Sigma\succ0$ cố định, nên có thể thay
$\vmu=\xbar$ trước rồi tối ưu theo $\Sigma$.
\nguon{Chương 4, §4.3, phân rã (4-14), tr.~169; Result 4.11, tr.~171--172}

\buoc{3}{Đổi biến để tìm ma trận hiệp phương sai}
Đặt $\Theta=\Sigma^{-1}$, gọi là \emph{ma trận độ chính xác}. Đây là
phép đổi biến một-một trên các ma trận xác định dương. Vì
$\log|\Sigma|=-\log|\Theta|$, log-hợp lý sau khi chọn trung bình là
\begin{equation}
 \ell(\xbar,\Theta)=K+\frac n2\log|\Theta|-\frac12\tr(\Theta A).
 \label{eq:mleprecision}
\end{equation}
Ý nghĩa của hai hạng: $\log|\Theta|$ thưởng cho độ phân tán nhỏ, còn
$\tr(\Theta A)$ phạt nếu độ phân tán quá nhỏ so với các độ lệch thực tế.
Cực đại cân bằng hai hạng đó.

Để đạo hàm, dùng hai quy tắc ma trận
\[
 d\log|\Theta|=\tr(\Theta^{-1}d\Theta),\qquad
 d\tr(\Theta A)=\tr(A\,d\Theta).
\]
Ký hiệu $d\Theta$ là một thay đổi nhỏ của ma trận, không phải vectơ
tịnh tiến $\mathbf d$ dùng ở phần biến đổi dữ liệu. Hai quy tắc trên là
bản mở rộng của $(\log\theta)'=1/\theta$ và $(a\theta)'=a$.
Từ \eqref{eq:mleprecision},
\[
 d\ell=\frac12\tr\bigl[(n\Theta^{-1}-A)d\Theta\bigr].
\]
Tại điểm dừng, biểu thức phải bằng 0 theo mọi hướng thay đổi đối xứng,
nên $n\Theta^{-1}-A=0$. Suy ra
\begin{equation}
 \widehat\Sigma=\widehat\Theta^{-1}=\frac An
 =\frac1n\sum_j(\mathbf x_j-\xbar)(\mathbf x_j-\xbar)'
 =\frac{n-1}{n}S.
 \label{eq:mlecov}
\end{equation}
Đó là nguồn gốc của mẫu số $n$; ta không chọn mẫu số theo ý thích.

\buoc{4}{Kiểm tra điểm dừng thực sự là cực đại}
Đạo hàm bằng 0 chưa đủ để kết luận. Với hướng đối xứng $H\ne0$, đạo
hàm bậc hai theo đường $\Theta+tH$ là
\[
 \frac{d^2}{dt^2}\ell(\xbar,\Theta+tH)\bigg|_{t=0}
 =-\frac n2\tr(\Theta^{-1}H\Theta^{-1}H)<0.
\]
Thật vậy, trace trong vế phải bằng tổng bình phương các trị riêng của
ma trận đối xứng $\Theta^{-1/2}H\Theta^{-1/2}$, nên dương khi $H\ne0$.
Log-hợp lý là hàm lõm nghiêm ngặt theo $\Theta$ trên miền xác định dương,
vì vậy điểm dừng tìm được là cực đại duy nhất. Kết hợp Bước 2, đây cũng
là cực đại đồng thời theo cả trung bình và hiệp phương sai.
\nguon{Kết quả MLE: Chương 4, §4.3, Result 4.11, tr.~171--172.
 Sách chứng minh bằng Result 4.10; ở đây dùng dẫn xuất tương đương theo ma trận độ chính xác}

\hieu{Trung bình mẫu đặt tâm tốt nhất của đám mây. Hiệp phương sai MLE
khớp độ phân tán quanh tâm đó bằng trung bình của các tích ngoài.
Tối ưu theo ma trận độ chính xác giúp nhìn rõ nghiệm và kiểm tra cực đại.}

\subsection{MLE khác ước lượng không chệch như thế nào?}
Thay các giá trị quan sát bằng biến ngẫu nhiên, các công thức MLE trở
thành các \emph{ước lượng} $\widehat{\vmu}=\Xbar$ và
$\widehat\Sigma=(n-1)S/n$. Khi đã thay số, đó là các \emph{giá trị ước lượng}.
Ta có
\[
 \E(S)=\Sigma,\qquad
 \E(\widehat\Sigma)=\frac{n-1}{n}\Sigma.
\]
$S$ không chệch; MLE hiệp phương sai bị chệch xuống ở mẫu hữu hạn.
MLE tối ưu hàm hợp lý, không tối ưu tiêu chí ``không chệch''. Hai mục
tiêu khác nhau, nên cho mẫu số khác nhau. Khi $n$ lớn, tỉ lệ $(n-1)/n$
tiến về 1 và khác biệt giảm.

\begin{example}[Tính MLE từ mẫu bốn quan sát]
Với $\xbar=(6,10)'$, $S$ trong \eqref{eq:sample4}, $n=4$,
\[
 A=3S=\begin{pmatrix}24&-10\\-10&6\end{pmatrix},\quad
 \widehat{\vmu}=\begin{pmatrix}6\\10\end{pmatrix},\quad
 \widehat\Sigma=\frac A4=\begin{pmatrix}6&-2.5\\-2.5&1.5\end{pmatrix}.
\]
Ta có $|\widehat\Sigma|=2.75>0$ và $\widehat\Theta=4A^{-1}$,
kiểm tra được $4\widehat\Theta^{-1}=A$.
Để kiểm định Hotelling, ta dùng $S$, không thế $\widehat\Sigma$ vào
công thức $n\mathbf r'S^{-1}\mathbf r$ mà giữ nguyên ngưỡng $F$.
Để thiết lập tỷ số hợp lý, ta dùng đúng MLE $\widehat\Sigma$.
\end{example}

\luuy{Nếu $A$ suy biến, $A/n$ không thuộc miền $\Sigma\succ0$.
Không có MLE hữu hạn trong mô hình chuẩn không suy biến này: có thể
co phương sai về 0 theo một hướng không có độ lệch quan sát, khiến hàm
hợp lý tăng không giới hạn. Tình huống này xảy ra khi $n\le p$, hoặc
do các cột phụ thuộc tuyến tính. Phụ lục~\ref{app:mle_global} giải thích
bằng trị riêng, đồng thời kiểm tra cực đại mà không dùng đạo hàm ma trận.}

\kiemtra{Dùng cùng một mẫu, $S$ và MLE hiệp phương sai có thể khác nhau
không? Có; MLE bằng $(n-1)S/n$. MLE trung bình có cần biết trước
$\Sigma$ không? Không, nghiệm $\xbar$ đúng với mọi $\Sigma\succ0$.}


\subsection{Không chuẩn nhưng mẫu lớn: luật số lớn và CLT}
Giả sử mẫu iid từ tổng thể bất kỳ, có trung bình $\vmu$ và hiệp phương
sai hữu hạn $\Sigma$, với $p$ cố định. Khi $n\to\infty$,
\begin{equation}
 \Xbar\xrightarrow{P}\vmu,\qquad S\xrightarrow{P}\Sigma.
 \label{eq:lln}
\end{equation}
Ký hiệu $\xrightarrow{P}$ là hội tụ theo xác suất: xác suất sai lệch lớn
hơn bất kỳ mức dương cho trước tiến về 0.
\nguon{Chương 4, §4.5, Result 4.12 và công thức (4-26), (4-27), tr.~175}

Định lý giới hạn trung tâm nhiều chiều nói rằng
\begin{equation}
 \sqrt n(\Xbar-\vmu)\xrightarrow{d}\N_p(0,\Sigma).
 \label{eq:clt}
\end{equation}
Ký hiệu $\xrightarrow{d}$ là hội tụ về phân phối. Công thức này mô tả
hình dạng giới hạn của sai lệch đã nhân $\sqrt n$; không nói mỗi quan sát
$\mathbf X_j$ trở thành chuẩn khi tăng $n$.
\nguon{Chương 4, §4.5, Result 4.13, tr.~176}

Nếu $\Sigma$ xác định dương, Slutsky cho phép thay $\Sigma$ bằng ước lượng
nhất quán $S$ trong dạng toàn phương:
\begin{equation}
 n(\Xbar-\vmu)'S^{-1}(\Xbar-\vmu)\xrightarrow{d}\chi^2_p.
 \label{eq:slutsky}
\end{equation}
\nguon{Chương 4, §4.5, công thức (4-28), tr.~176; ứng dụng ở Result 5.4, tr.~235}

\hieu{CLT cho hình dạng chuẩn ở giới hạn, luật số lớn cho phép ước lượng
độ phân tán, và Slutsky nối hai kết quả. Không cần giả thiết chuẩn tổng thể
cho kết quả giới hạn này, nhưng cần iid, moment hữu hạn và hiệp phương sai
không suy biến.}

Trong mẫu chuẩn, các kết quả Wishart và độc lập là \emph{chính xác ở mẫu
hữu hạn}; trong tổng thể bất kỳ, chúng không tự được khôi phục chỉ vì mẫu
lớn. Ta chỉ dùng kết quả xấp xỉ thích hợp. Cách chọn sẽ được tóm tắt ở
Phần~\ref{sec:large}.

\kiemtra{Có cần $n-p\ge30$ để $S\xrightarrow{P}\Sigma$ đúng không?
Không. Hội tụ là phát biểu khi $n\to\infty$; một ngưỡng mẫu hữu hạn không
phải giả thiết của luật số lớn.}
